% Moved to supplement from Results VII-E (EVT sensitivity + failure decomposition).
\section{Extended EVT Sensitivity and Failure-Mechanism Decomposition}
\label{sup:evt_detail}

\begin{table}[t]
\centering
\small
\caption{Sensitivity Analysis of Extreme Value Theory (SPOT) Calibration Parameters across Benchmark Streams. We evaluate the sensitivity of empirical False Positive Rates (on held-out nominal validation splits and test sequences) and Point-F1 across initial tail quantiles $u \in \{0.95, 0.98, 0.99\}$, design risk parameters $q \in \{10^{-2}, 10^{-3}, 10^{-4}\}$, and kurtosis adaptation coefficients $\beta \in \{0.0, 0.083, 0.167, 0.25, 0.50\}$ alongside non-parametric baselines.}
\label{tab:evt_sensitivity}
\resizebox{\linewidth}{!}{%
\begin{tabular}{@{}l cccc c@{}}
\toprule
\textbf{Threshold Calibration Method} & \textbf{Design $q$} & \textbf{Tail $u$} & \textbf{Val FPR} & \textbf{Test FPR} & \textbf{Point-F1} \\
\midrule
\multicolumn{6}{l}{\textit{\textbf{Initial Tail Quantile and Risk Level Sensitivity (Static SPOT)}}} \\
Static SPOT ($u=0.95, q=10^{-4}$)      & $10^{-4}$ & 0.95 & 0.00009 & 0.1727 & 0.1029 \\
Static SPOT ($u=0.95, q=10^{-3}$)      & $10^{-3}$ & 0.95 & 0.00190 & 0.1762 & 0.1030 \\
Static SPOT ($u=0.95, q=10^{-2}$)      & $10^{-2}$ & 0.95 & 0.00741 & 0.1895 & 0.1020 \\
\addlinespace[2pt]
Static SPOT ($u=0.98, q=10^{-4}$)      & $10^{-4}$ & 0.98 & 0.00009 & 0.1733 & 0.1028 \\
Static SPOT ($u=0.98, q=10^{-3}$) (Baseline)& $10^{-3}$ & 0.98 & 0.00190 & 0.1764 & 0.1030 \\
Static SPOT ($u=0.98, q=10^{-2}$)      & $10^{-2}$ & 0.98 & 0.00743 & 0.1886 & 0.1022 \\
\addlinespace[2pt]
Static SPOT ($u=0.99, q=10^{-4}$)      & $10^{-4}$ & 0.99 & 0.00017 & 0.1736 & 0.1030 \\
Static SPOT ($u=0.99, q=10^{-3}$)      & $10^{-3}$ & 0.99 & 0.00210 & 0.1763 & 0.1029 \\
Static SPOT ($u=0.99, q=10^{-2}$)      & $10^{-2}$ & 0.99 & 0.00703 & 0.1862 & 0.1029 \\
\midrule
\multicolumn{6}{l}{\textit{\textbf{Kurtosis-Adaptive Multiplier Sensitivity ($\tau_{\text{eff}} = \tau(1 + \beta\tanh\gamma_2)$)}}} \\
Kurtosis-Adaptive ($\beta=0.000$, Standard) & $10^{-3}$ & 0.98 & 0.00190 & 0.1764 & 0.1030 \\
Kurtosis-Adaptive ($\beta=0.083$)           & $10^{-3}$ & 0.98 & 0.27780 & 0.2241 & 0.0973 \\
Kurtosis-Adaptive ($\beta=0.167$)           & $10^{-3}$ & 0.98 & 0.30463 & 0.2686 & 0.0964 \\
Kurtosis-Adaptive ($\beta=0.250$)           & $10^{-3}$ & 0.98 & 0.33991 & 0.3813 & 0.0995 \\
Kurtosis-Adaptive ($\beta=0.500$)           & $10^{-3}$ & 0.98 & 0.54032 & 0.6096 & 0.2560 \\
Moors Quantile Kurtosis ($\beta=0.167$)     & $10^{-3}$ & 0.98 & 0.19350 & 0.1848 & 0.1024 \\
\midrule
\multicolumn{6}{l}{\textit{\textbf{Non-Parametric Validation Baselines}}} \\
Max-of-Validation                           & --        & --   & \textbf{0.00000} & \textbf{0.1774} & 0.1029 \\
Empirical 99.5th Percentile                 & --        & --   & 0.00544 & 0.1837 & 0.1028 \\
Gaussian 3-Sigma ($\mu + 3\sigma$)          & --        & --   & 0.00882 & 0.1879 & 0.1016 \\
\bottomrule
\end{tabular}%
}
\end{table}


Table~\ref{tab:evt_sensitivity} provides a multi-parameter sensitivity analysis of the extreme-value calibration procedure across varying initial tail thresholds $u$, theoretical design risk levels $q$, and kurtosis-adaptive coefficients $\beta$. Three key findings emerge:
First, varying the initial threshold from $u=0.95$ to $u=0.99$ produces negligible changes in out-of-sample test FPR ($17.27\%\text{--}17.36\%$) and Point-F1 ($0.1028\text{--}0.1030$), indicating that GPD tail estimates are stable with respect to initial tail percentile selection.
Second, scaling the theoretical design risk $q$ by two orders of magnitude (from $10^{-2}$ to $10^{-4}$) shifts validation FPR from $0.0074$ to $0.00009$, but test FPR shifts by less than $1.7\%$ (from $18.95\%$ to $17.27\%$), demonstrating that operational distribution drift completely dominates the theoretical design parameter.
Third, increasing the kurtosis adjustment coefficient $\beta$ systematically degrades performance: at $\beta=0.083$, held-out validation FPR jumps from $0.0019$ to $0.2778$, and at $\beta=0.50$, it collapses to $0.5403$ (test FPR $60.96\%$), flooding the system with alarms. Non-parametric validation rules (such as Max-of-Validation and Empirical 99.5th Percentile) achieve comparable or superior nominal calibration ($0.0000\text{--}0.0054$ validation FPR) without requiring parametric tail assumptions.

\begin{table*}[t]
\centering
\caption{Empirical Decomposition of Extreme Value Calibration Failure. We isolate the four hypothesized drivers of threshold breakdown across benchmark telemetry streams: (1) residual autocorrelation ($\rho_1$) and extremal clustering ($\theta, 1/\theta$), which reduce the effective independent tail sample size ($N_{\text{eff}} = \theta N_u$); (2) non-stationary distribution drift between calibration and held-out evaluation splits (Kolmogorov-Smirnov $D_{\text{KS}}$ and relative 98th-percentile drift $\Delta q_{98}\%$); and (3) actual held-out nominal False-Positive Rate and empirical risk overshoot relative to the theoretical design risk $q=10^{-3}$ ($0.1\%$). Observation-space models exhibit large distribution drift ($\Delta q_{98} = +470\%$) and severe clustering ($1/\theta \approx 12.7\text{--}13.3$ steps), causing a $43\times$ to $47\times$ overshoot of design risk. In contrast, latent prediction compresses distribution drift ($\Delta q_{98} = -0.36\%$), reducing risk overshoot by an order of magnitude.}
\label{tab:evt_failure_decomposition}
\small
\begin{tabular}{lcccccc} \toprule
\textbf{Architecture \& Objective} & \textbf{Autocorr. ($\rho_1$)} & \textbf{Extremal Index ($\theta$)} & \textbf{Mean Cluster} & \textbf{KS Drift ($D_{\text{KS}}$)} & \textbf{Held-Out FPR} & \textbf{Risk Overshoot} \\
 & (Lag-1) & (Ferro-Segers) & ($1/\theta$ steps) & (Cal vs.\ Eval) & (Nominal Target: $0.1\%$) & ($\text{FPR}_{\text{eval}} / 10^{-3}$) \\ \midrule
TimesNet (Reconstructive) & $0.983 \pm 0.014$ & $0.528 \pm 0.518$ & $12.7 \pm 17.1$ & $0.357 \pm 0.178$ & $4.32\% \pm 7.60\%$ & $43.2\times \pm 76.0\times$ \\
TranAD (Reconstructive) & $0.990 \pm 0.007$ & $0.484 \pm 0.483$ & $13.3 \pm 17.1$ & $0.388 \pm 0.284$ & $4.76\% \pm 10.43\%$ & $47.6\times \pm 104.3\times$ \\
TS-JEPA (Latent Predictive) & $0.820 \pm 0.135$ & $0.483 \pm 0.356$ & $5.2 \pm 5.8$ & $0.154 \pm 0.069$ & $0.33\% \pm 0.59\%$ & $3.3\times \pm 5.9\times$ \\
\bottomrule
\end{tabular}
\end{table*}


To diagnose why extreme-value calibration breaks down on operational telemetry, Table~\ref{tab:evt_failure_decomposition} isolates four statistical mechanisms across benchmark telemetry streams:
\begin{enumerate}
    \item \textbf{Residual Autocorrelation and Extremal Clustering:} The Pickands-Balkema-de Haan theorem \cite{balkema1974residual} requires mutually independent tail exceedances. However, reconstructive observation residuals exhibit serial correlation, with lag-1 autocorrelation $\rho_1 = 0.983 \pm 0.014$ for TimesNet and $\rho_1 = 0.990 \pm 0.007$ for TranAD. Using the Ferro-Segers extremal index estimator \cite{ferro2003inference}, observation exceedances form persistent temporal clusters with mean cluster sizes of $1/\theta = 12.7$ to $13.3$ steps (reaching $45.0$ steps on GECCO water quality telemetry). Clustering collapses the effective independent tail sample size ($N_{\text{eff}} = \theta N_u$) from $N_u = 77\text{--}91$ observations down to $N_{\text{eff}} \approx 2\text{--}6$ independent clusters. In contrast, latent prediction compresses lag-1 autocorrelation to $\rho_1 = 0.820 \pm 0.135$ and mean cluster size to $5.2 \pm 5.8$ steps.
    \item \textbf{Non-Stationary Distribution Drift:} Non-stationarity drives a wedge between calibration and held-out evaluation splits. Two-sample Kolmogorov-Smirnov tests reveal substantial drift in reconstructive models ($D_{\text{KS}} = 0.357 \pm 0.178$ for TimesNet; $D_{\text{KS}} = 0.388 \pm 0.284$ for TranAD; $p < 10^{-13}$ across streams). This drift shifts the 98th percentile of nominal residuals by $+470.8\%$ on TimesNet and $+48.6\%$ on TranAD between splits. Conversely, TS-JEPA maps temporal sequences into an abstract representation space where environmental shifts are compressed, yielding less than half the KS distance ($D_{\text{KS}} = 0.154 \pm 0.069$) and near-zero relative tail drift ($\Delta q_{98} = -0.36\% \pm 14.1\%$).
    \item \textbf{GPD Goodness-of-Fit Degradation:} Cram{\'e}r-von Mises tests for Generalized Pareto tail convergence reject the stationary GPD null hypothesis ($p < 0.05$) across autocorrelated observation streams. Because parameter estimation fits a stationary tail model to non-stationary, autocorrelated bursts, parametric SPOT thresholds provide no advantage over simple empirical validation percentiles.
    \item \textbf{Empirical Risk Overshoot:} Against the specified design risk $q = 10^{-3}$ ($0.1\%$), TimesNet generates a held-out nominal FPR of $4.32\% \pm 7.60\%$ ($43.2\times$ overshoot), and TranAD produces $4.76\% \pm 10.43\%$ ($47.6\times$ overshoot). By suppressing residual drift and tail clustering, TS-JEPA curbs the empirical nominal FPR to $0.33\% \pm 0.59\%$ ($3.3\times$ overshoot), reducing the discrepancy between theoretical risk and operational deployment by an order of magnitude.
\end{enumerate}

